% !PRINTED 721
\weberpage{721}

The number $1 + \mu$ has the indices $\alpha = 0$, $\beta = \nu'$, and
consequently, by (7),
\[
\chi(1+\mu) = \Theta^{\nu'} = -1,
\]
if $\chi$ is a primitive character.

Further, for an odd $n$:
\[
n(1+\mu) \equiv n + \mu \pmod{m},
\]
and consequently
\[
\chi(n+\mu) = -\chi(n).
\]
If we now consider the function $f(r^s)$:
\weq{10}{f(r^s) = \overset{n}{\Sigma} \chi(n)\, r^{sn},}
in which $n$ runs through the odd numbers of a complete system of
residues for the modulus $m$, then, since in (10) we may replace $n$ by
$n+\mu$ under the summation sign, we obtain:
\[
f(r^s) = -r^{\mu s} \sum \chi(n)\, r^{sn},
\]
and since $r^\mu = -1$:
\[
f(r^s) = -(-1)^s f(r^s).
\]
Thus
\weq{11}{f(r^s) = 0, \quad \text{\gesperrt{if $s$ is even,}}}
\gesperrt{and accordingly we may, in the formulae (9), restrict the summation
letter $s$ to the odd numbers that are less than $m$.}

But by (10) we have:
\[
f(r^s) = \chi(s)^{-1} \overset{n}{\Sigma} \chi(ns)\, r^{sn},
\]
and if one replaces $sn$ by $n$, which is permitted when $s$ is odd:
\weq{12}{f(r^s) = \chi(s)^{-1} f(r), \quad \text{\gesperrt{if $s$ is odd.}}}
The function $f(r)$ has, when $\alpha, \beta$ are the indices of $n$,
the expression
\weq{13}{f(r) = \overset{n}{\Sigma} (\pm 1)^\alpha \Theta^\beta r^n,}
and is therefore identical in meaning with the cyclotomic resolvent
considered in \S.\,17 of this volume
\[
(\pm 1, \Theta, r)
\]
for which the relation was set up in \S.\,17, (17):
\weq{14}{(\pm 1, \Theta, r)(\pm 1, \Theta^{-1}, r) = \pm m.}
Here the upper signs hold throughout for the sums $X_1$, the lower for
the sums $X_2$ [\S.\,195, (7)].

We now divide the sequence of the numbers $s$ occurring in (9) into
four parts. Let $t$ be a symbol which
%CONT>>
